% ============================================================
% 分节契约
% purpose:      建立"为什么这个问题重要、为什么现有方案不够、本文做了什么"
% allowed claims: 版图性陈述须有文献行支撑(01-literature 精读卡)
% forbidden claims: 贬低竞品的无证据表述;防御性让步(anti-defensive-writing)
% inputs:       01-topic.md · 00-canonical/evidence_table.md
% 验证方式:     nature-citation 配引用 + ref-verifier 核对
% ============================================================
\section{Introduction}\label{sec:intro}

% 叙事顺序按 anti-defensive-writing 发布会原则:
% 第 1 段问题与利害 → 第 2 段现状与缺口(引文)→ 第 3 段本文做法
% → 第 4 段贡献清单 → 第 5 段结果预告
<第 1 段:领域问题与利害。>

<第 2 段:现有方法路线与缺口,配 \cite{...}。>

<第 3 段:本文方法核心思想。>

Our contributions are as follows:
\begin{itemize}
  \item <贡献 1:与 argument_map 的主 claim 一一对应>
  \item <贡献 2>
  \item <贡献 3>
\end{itemize}
